\section{Introduction}
\label{sec:intro}

A government that guarantees mortgages, lends for education and funds itself from payroll and
income taxes has taken a position on wages without ever describing it as one. Nothing in the
public accounts says so, because the accounts were not built to say it. They record what the
government owes and what it owns. They do not record which income has to keep flowing for the
claims it stands behind to be paid.

That gap is not peculiar to government. The financial accounts of the United States describe the
whole financial system as a matrix of claims between sectors, and for seventy years they have
answered two questions well: who owes this claim, and who holds it. A third question sits
underneath both and is not asked. A mortgage is paid out of household wages. A corporate bond is
paid out of business revenue. Treasury debt is paid out of federal receipts, which are themselves
drawn in measurable proportions from wages and from elsewhere. Classify every claim by the income
that services it, and the claim stock sorts into groups that rise and fall together for reasons
that have nothing to do with who happens to hold the paper.

This paper builds that classification and crosses it with the holder structure the accounts
already carry. We call the result labor backing accounts.

\paragraph{Why the question matters to people, not only to statisticians.} A statistic that
locates wage dependence is a statistic about who absorbs a fall in pay. The household that owes
a car loan and a student loan out of a single wage does not experience a labor share as an
aggregate. It experiences a payment that does not move when the income behind it does. Our
measurements say that such households are not a minority at the edge of the system: claims
attributed to labor income are the majority of US debt, and they are spread far more evenly
across the wage distribution than the wages attributed to servicing them. So the ratio of
claims to wages rises as pay falls. That is a property of the claim structure rather than a
prediction about anybody's behavior, and it is a statement about stocks: whether a
proportional fall in pay is regressive in welfare, in default or in debt-service terms needs
payment schedules, maturity, transfers and expenditure, which these accounts do not carry.

The same measurement speaks to who ends up carrying the loss. On an agency-guaranteed mortgage
the credit risk sits with the guarantor, not with the pension fund holding the security. On
Treasury debt serviced from wage taxes the exposure is a revenue shortfall on a claim the
government owes. Counting only what the public sector holds understates its position by more
than a factor of two. Whether that concentration is prudent is a question for others. Whether it
exists, and how large it is, is a question of measurement, and it has not been answered.

\paragraph{Research questions and objectives.} The paper is organized around three questions,
each paired with an objective and the section that meets it.

\begin{description}[leftmargin=2.2em,style=nextline]
  \item[RQ1] Which income services the stock of US financial claims, and how much of it is
  wages? \\
  \textbf{RO1.} Define labor backing at the level of a claim class, measure the coefficient for
  each of thirteen classes from household surveys and national accounts, and report the resulting
  ratio under every convention we can defend. Section~\ref{sec:accounts}.
  \item[RQ2] Who holds, guarantees or owes the claims that wages pay, and has that changed? \\
  \textbf{RO2.} Build the holder map of the wage-backed stock, state the consolidation rule that
  decides what it means, and run the accounts back to 1947 so the present level can be read
  against its own history. Section~\ref{sec:accounts}.
  \item[RQ3] How is that stock distributed across the households whose wages service it? \\
  \textbf{RO3.} Split the household part of the stock by wage quintile and set it beside the
  liquid buffers of the same households. Section~\ref{sec:households}.
\end{description}

\paragraph{Contributions.} The paper makes three.

\begin{enumerate}[leftmargin=1.6em]
  \item \textbf{A measurement framework.} Labor backing accounts decompose the stock of US
  financial claims by servicing income and by holder at the same time, on the same claim stock in
  the same period. The framework yields a ratio, a federal exposure share, a holder map for both
  sides, a series from 1952 and a breakdown by wage quintile, with every convention and its
  sensitivity stated in full.
  \item \textbf{A holder map of two differently observable sides.} The wage side is a measured
  stock. The side that would replace it, claims on AI capital, is assembled from company filings
  and is therefore reported as bounds rather than as a point. The two are placed beside each
  other and never netted, because netting would require a dollar base the second side does not
  have.
  \item \textbf{A distributional reading, and an honest null.} Wage-backed claims per dollar of
  wages fall as pay rises, so the claims are spread more evenly than the income servicing them.
  Against that, a pre-registered test of whether the measure predicts bank credit losses returned
  nothing, twice and on two samples. We report it in Section~\ref{sec:limitations} because it
  bounds what the accounts are for.
\end{enumerate}

\paragraph{What this paper does not claim.} We do not forecast automation, and we do not claim
the fiscal mechanism, which is established in the work reviewed in Section~\ref{sec:related}. We
do not claim the framework is new in the strong sense: our search located no prior construction
of the same object, but that search was not systematic, so the claim is ``none located'' rather
than ``none exists''. One priority claim in this line of work was retired after a non-systematic
search found prior work on it, and the same could happen here. Above all, these accounts are
accounting. They say where wage dependence sits. They do not say which institution will take a
loss, and Section~\ref{sec:limitations} reports the test that establishes it.
